\documentclass[12pt]{article}
\usepackage[a4paper,margin=1in]{geometry}
\usepackage{amsmath,amssymb}
\usepackage{enumitem}
\usepackage[a4paper,margin=1in]{geometry}
\usepackage{amsmath,amssymb}
\usepackage{hyperref}
\usepackage{titlesec}
\title{Main results}
\author{}
\date{}
\begin{document}
\maketitle
\vspace{-2em}

\section*{Cesàro Summation}

A series $\sum a_n$ with partial sums $S_N=\sum_{k=0}^{N}a_k$ is \emph{Cesàro summable} (C,1) to $L$ if  
\[
\sigma_N=\frac{1}{N+1}\sum_{n=0}^{N}S_n \;\longrightarrow\; L.
\]
Every convergent series is Cesàro summable to the same limit, but not conversely.  

\section*{Positive Summation Kernels}
A family of functions $(K_n)_{n\ge1}$ on $\mathbb{R}$ (or on $\mathbb{T}$, the circle) is called a \emph{positive summation kernel} if the following hold:
\begin{enumerate}[label=(\roman*)]
  \item $K_n(x) \ge 0$ for all $x \in \mathbb{R}$ and each $n$;
  \item $\displaystyle \int_{-\infty}^{\infty} K_n(x)\,dx = 1$ for all $n$;
  \item for every $\delta > 0$, 
  \[
    \int_{|x|\ge \delta} K_n(x)\,dx \;\longrightarrow\; 0 \quad \text{as } n \to \infty.
  \]
\end{enumerate}
Equivalently, $(K_n)$ is an \emph{approximate identity}: the mass of $K_n$ concentrates near $x=0$ as $n\to\infty$.  
For such kernels, the convolution $(f*K_n)(x) = \int f(x-t)K_n(t)\,dt$ satisfies
\[
  (f*K_n)(x) \to f(x)
\]
for all $x$ where $f$ is continuous.

\section*{Riemann-Lebesgue Lemma}

\noindent Theorem 2.2 (Riemann--Lebesgue lemma)
Let $f$ be absolutely integrable in the Riemann sense on a finite or infinite interval $I$. Then
\[
  \lim_{\lambda\to\infty} \int_{I} f(u)\sin(\lambda u)\,du = 0.
\]
(A corresponding result holds for $\cos(\lambda u)$ or $e^{i\lambda u}$.)

\subsection*{Corollary for Fourier Coefficients}

\noindent Corollary.
Let $f \in L^1(\mathbb{T})$ (i.e., $\int_{-\pi}^{\pi} |f(x)| dx < \infty$). Its complex Fourier coefficients are:
\[
  \hat{f}(n) = c_n = \frac{1}{2\pi}\int_{-\pi}^{\pi} f(x)e^{-inx}\,dx, \qquad n\in\mathbb{Z}.
\]
\emph{The Fourier coefficients of any integrable function vanish at infinity:}
\[
  \hat{f}(n) \longrightarrow 0 \quad \text{as } |n| \to \infty.
\]

Proof Sketch of Corollary:
The Fourier coefficient $\hat{f}(n)$ can be expressed as:
$$\hat{f}(n) = \frac{1}{2\pi} \left( \int_{-\pi}^{\pi} f(x) \cos(nx) \, dx + i \int_{-\pi}^{\pi} f(x) \sin(-nx) \, dx \right)$$
Since $\cos(nx)$ and $\sin(nx)$ are bounded functions, the absolute integrability of $f$ is preserved when multiplied by these functions. As $|n| \to \infty$, both the real and imaginary parts of the integral tend to zero by applying the Riemann--Lebesgue lemma (and its cosine variant) with $\lambda = n$. Therefore, $\hat{f}(n) \to 0$ as $|n| \to \infty$.

\section*{Convolution (on $\mathbb{R}_+$) and the Convolution Theorem}

\noindent
Assume $f,g:\mathbb{R}_+\to\mathbb{C}$ are causal (zero for $t<0$), piecewise continuous, and of exponential order so that Laplace transforms exist.

\subsection*{Definition (Convolution)}
\[
(f*g)(t)\;:=\;\int_{0}^{t} f(t-\tau)\,g(\tau)\,d\tau,\qquad t\ge 0.
\]

\subsection*{Algebraic Laws}
For all admissible $f,g,h$ and scalars $\alpha,\beta$:
\begin{align*}
\text{Commutative: }& f*g = g*f,\\
\text{Associative: }& (f*g)*h = f*(g*h),\\
\text{Distributive: }& f*(\alpha g+\beta h) = \alpha(f*g)+\beta(f*h).
\end{align*}

\subsection*{Convolution Theorem (Laplace)}
Let $F(s)=\mathcal{L}\{f\}(s)$ and $G(s)=\mathcal{L}\{g\}(s)$ on a common half-plane $\Re s>\sigma_0$. Then
\[
\;\mathcal{L}\{f*g\}(s)=F(s)\,G(s)\;
\]
(Convergence holds on any half-plane where both $F$ and $G$ exist.)
\section*{The Laplace Transform}

Let $H(t)$ be the Heaviside step function, $H(t)=0$ for $t<0$ and $H(t)=1$ for $t\ge0$.  
We will always work with \emph{causal} functions, i.e. we take $f$ to be defined as
\[
f(t)\;:=\;f(t)\,H(t),
\]
so that $f(t)=0$ for all $t<0$.

\medskip
\noindent For a complex parameter $s$ with $\Re s$ large enough so that the integral converges (the \emph{region of convergence}),
the \emph{Laplace transform} of $f$ is
\[
\mathcal{L}\{f\}(s)\;=\;\int_{0}^{\infty} f(t)\,e^{-st}\,dt.
\]
We also write $\tilde f(s)$ or $F(s)$ for $\mathcal{L}\{f\}(s)$.

\subsection*{Functions of Exponential Order}

\noindent
Let $k>0$. A function $f:[0,\infty)\to\mathbb{C}$ is said to be of \emph{exponential order $k$} if:

\begin{enumerate}[label=(\roman*)]
  \item $f$ is continuous on every finite subinterval of $[0,\infty)$, except possibly for finitely many jump discontinuities;
  \item there exists a constant $M>0$ such that
  \[
    |f(t)| \le M e^{k t}, \qquad t \ge 0.
  \]
\end{enumerate}

\noindent
The class of all such functions is denoted by $\mathcal{E}_k$.  
If $f \in \mathcal{E}_k$ for some $k$, we say that $f$ is of \emph{exponential type} and write $f \in \mathcal{E} = \bigcup_{k>0} \mathcal{E}_k$.

\medskip
\noindent
\textbf{Theorem.}  
If $f \in \mathcal{E}_k$, then the Laplace transform
\[
\tilde{f}(s) = \int_{0}^{\infty} f(t)e^{-st}\,dt
\]
exists for all $s > k$.

\subsection*{Common Laplace Transforms}
(Here $a,b\in\mathbb{R}$, $n\in\mathbb{N}$, and the stated condition on $s$ ensures convergence.)
\\
\\
\\
\noindent
\begin{tabular*}{\textwidth}{@{\extracolsep{\fill}} l l l @{}}
\textbf{Function $f(t)$} & \textbf{$\mathcal{L}\{f\}(s)$} & \textbf{Condition on $s$}\\[4pt]\hline
$1$ & $\displaystyle \frac{1}{s}$ & $\Re s>0$\\[8pt]
$t^n$ & $\displaystyle \frac{n!}{s^{\,n+1}}$ & $\Re s>0$\\[8pt]
$e^{at}$ & $\displaystyle \frac{1}{s-a}$ & $\Re s>a$\\[8pt]
$\sin(bt)$ & $\displaystyle \frac{b}{s^{2}+b^{2}}$ & $\Re s>0$\\[8pt]
$\cos(bt)$ & $\displaystyle \frac{s}{s^{2}+b^{2}}$ & $\Re s>0$\\[8pt]
$\sinh(bt)$ & $\displaystyle \frac{b}{s^{2}-b^{2}}$ & $\Re s>|b|$\\[8pt]
$\cosh(bt)$ & $\displaystyle \frac{s}{s^{2}-b^{2}}$ & $\Re s>|b|$\\[8pt]
$f(t-a)\,H(t-a)$ & $\displaystyle e^{-as}F(s)$ & $a>0$ (second shifting)\\[8pt]
$e^{at}f(t)$ & $\displaystyle F(s-a)$ & $\Re s>a+\sigma_0$, if $F$ converges for $\Re s>\sigma_0$
\end{tabular*}

\subsection*{Operational Rules for the Laplace Transform}

Let $\mathcal{L}\{f(t)\}(s) = \tilde{f}(s)$, and assume all transforms exist for sufficiently large $\Re s$.

\begin{enumerate}[label=(\roman*)]
  \item \textbf{Linearity:}
  \[
  \mathcal{L}\{\alpha f(t) + \beta g(t)\}(s)
  = \alpha \tilde{f}(s) + \beta \tilde{g}(s),
  \qquad \alpha, \beta \in \mathbb{C}.
  \]

  \item \textbf{Exponential (Damping) Rule:}
  \[
  \mathcal{L}\{e^{a t} f(t)\}(s)
  = \tilde{f}(s - a).
  \]

  \item \textbf{Time Shift (Delay) Rule:}
  If $f(t) = 0$ for $t < 0$ and $a > 0$, then
  \[
  \mathcal{L}\{f(t - a)\}(s)
  = e^{-a s}\,\tilde{f}(s).
  \]

  \item \textbf{Scaling Rule:}
  For $a > 0$,
  \[
  \mathcal{L}\{f(a t)\}(s)
  = \frac{1}{a}\,\tilde{f}\!\left(\frac{s}{a}\right).
  \]
  \item \textbf{Multiplication by $t$ (and by $t^n$):}
  \[
    \mathcal{L}\{t\,f(t)\}(s) \;=\; -\frac{d}{ds}\tilde f(s), 
    \qquad
    \mathcal{L}\{t^{n}f(t)\}(s) \;=\; (-1)^{n}\frac{d^{\,n}}{ds^{\,n}}\tilde f(s)\quad (n\in\mathbb{N}).
  \]

  \item \textbf{Derivatives in $t$:}
  \[
    \mathcal{L}\{f'(t)\}(s)=s\,\tilde f(s)-f(0^+),\qquad
    \mathcal{L}\{f''(t)\}(s)=s^{2}\tilde f(s)-s f(0^+)-f'(0^+),
  \]
  and in general
  \[
    \mathcal{L}\{f^{(n)}(t)\}(s)=s^{n}\tilde f(s)-\sum_{k=0}^{n-1}s^{\,n-1-k}f^{(k)}(0^+).
  \]

  \item \textbf{Integral in $t$:}
  \[
    \mathcal{L}\!\left\{\int_{0}^{t} f(u)\,du\right\}(s)=\frac{\tilde f(s)}{s}.
  \]
\end{enumerate}

\section*{Fourier Series Coefficients}

The representation of a $2\pi$-periodic function $f$ as a Fourier series depends on whether the complex exponential basis or the real trigonometric basis is used.

\subsection*{1. Complex Fourier Series}
For a function $f \in L^1(\mathbb{T})$, the series representation is:
\[
  f(x) \sim \sum_{n \in \mathbb{Z}} c_n e^{inx}
\]
The complex Fourier coefficients $c_n = \hat{f}(n)$ are computed as:
\[
  c_n = \hat{f}(n) = \frac{1}{2\pi}\int_{-\pi}^{\pi} f(x)e^{-inx}\,dx, \qquad n\in\mathbb{Z}.
\]

\subsection*{2. Real Fourier Series (Trigonometric)}
For a real-valued function $f$, the series representation is:
\[
  f(x) \sim a_0 + \sum_{n=1}^{\infty} \left( a_n \cos(nx) + b_n \sin(nx) \right)
\]
The real Fourier coefficients are computed as:
\begin{itemize}
    \item $a_0$ (Constant component):
    \[
      a_0 = \frac{1}{2\pi}\int_{-\pi}^{\pi} f(x)\,dx
    \]
    \item $a_n$ (Cosine coefficients, $n \ge 1$):
    \[
      a_n = \frac{1}{\pi}\int_{-\pi}^{\pi} f(x) \cos(nx)\,dx, \qquad n\in\mathbb{N}
    \]
    \item $b_n$ (Sine coefficients, $n \ge 1$):
    \[
      b_n = \frac{1}{\pi}\int_{-\pi}^{\pi} f(x) \sin(nx)\,dx, \qquad n\in\mathbb{N}
    \]
\end{itemize}
\emph{Note:} The relationship between the coefficients is $c_0 = a_0$, and for $n \ge 1$, $c_n = \frac{1}{2}(a_n - ib_n)$ and $c_{-n} = \frac{1}{2}(a_n + ib_n)$.

\section*{Fourier Inversion Theorem (Periodic Case)}

Let $f : \mathbb{R} \to \mathbb{C}$ be a $2\pi$-periodic function such that $f$ is piecewise continuous and of bounded variation on $[-\pi,\pi]$.  
Its (complex) Fourier coefficients are
\[
\hat{f}(n) = \frac{1}{2\pi}\int_{-\pi}^{\pi} f(t)e^{-int}\,dt, \qquad n \in \mathbb{Z}.
\]
Then the Fourier series
\[
S_N(f)(t) = \sum_{|n|\le N}\hat{f}(n)e^{int}
\]
converges, for each $t$, to
\[
\lim_{N\to\infty}S_N(f)(t)
= \frac{f(t^-)+f(t^+)}{2}.
\]
That is, the series converges to the midpoint of the left and right limits of $f$ at each point.  
In particular, if $f$ is continuous at $t$, then
\[
\sum_{n\in\mathbb{Z}}\hat{f}(n)e^{int} = f(t).
\]

\subsection*{Consequences}

\begin{itemize}
  \item If $\displaystyle \sum_{n\in\mathbb{Z}} |\hat{f}(n)| < \infty$, the Fourier series converges \emph{uniformly} to a \emph{continuous} $2\pi$-periodic function.
  \item If $f$ has a jump discontinuity at $t_0$, the Fourier series converges to $\tfrac{1}{2}(f(t_0^-)+f(t_0^+))$, not to either one-sided limit.
  \item Evaluating at $t=0$ gives
  \(
    \sum_{n\in\mathbb{Z}}\hat{f}(n)
    = \frac{f(0^-)+f(0^+)}{2},
  \)
\end{itemize}

\section*{Fourier Transform (on $\mathbb{R}$)}

\subsection*{Definition and Convention}
For $f \in L^1(\mathbb{R})$, the Fourier Transform $\hat{f}$ is defined as:
\[
\hat{f}(\omega) = \mathcal{F}\{f\}(\omega) = \int_{-\infty}^{\infty} f(t) e^{-i\omega t} dt, \qquad \omega \in \mathbb{R}.
\]

\subsection*{Real and Imaginary Parts}
For a real-valued function $f(t)$, the Fourier Transform $\hat{f}(\omega)$ is generally complex and can be split into real and imaginary parts:
\[
\hat{f}(\omega) = \int_{-\infty}^{\infty} f(t) (\cos(-\omega t) + i \sin(-\omega t)) dt = \int_{-\infty}^{\infty} f(t) \cos(\omega t) dt - i \int_{-\infty}^{\infty} f(t) \sin(\omega t) dt
\]

\noindent
Thus, we define:
\begin{align*}
\Re\{\hat{f}(\omega)\} &= \int_{-\infty}^{\infty} f(t) \cos(\omega t) dt \\
\Im\{\hat{f}(\omega)\} &= -\int_{-\infty}^{\infty} f(t) \sin(\omega t) dt
\end{align*}

\subsection*{Symmetry Simplifications (for Real-Valued $f$)}
The symmetry of $f(t)$ simplifies the computation of $\hat{f}(\omega)$.

\begin{itemize}[leftmargin=*, label={$\bullet$}]
    \item ($f(t) = f(-t)$):
    The product $f(t)\sin(\omega t)$ is odd, so $\int_{-\infty}^{\infty} f(t) \sin(\omega t) dt = 0$.
    The product $f(t)\cos(\omega t)$ is even, so $\int_{-\infty}^{\infty} f(t) \cos(\omega t) dt = 2 \int_{0}^{\infty} f(t) \cos(\omega t) dt$.
    \[
    \hat{f}(\omega) = 2 \int_{0}^{\infty} f(t) \cos(\omega t) dt \quad \text{($\hat{f}$ is purely real and even)}
    \]

    \item ($f(t) = -f(-t)$):
    The product $f(t)\cos(\omega t)$ is odd, so $\int_{-\infty}^{\infty} f(t) \cos(\omega t) dt = 0$.
    The product $f(t)\sin(\omega t)$ is even, so $\int_{-\infty}^{\infty} f(t) \sin(\omega t) dt = 2 \int_{0}^{\infty} f(t) \sin(\omega t) dt$.
    \[
    \hat{f}(\omega) = -i \left( 2 \int_{0}^{\infty} f(t) \sin(\omega t) dt \right) \quad \text{($\hat{f}$ is purely imaginary and odd)}
    \]
\end{itemize}

\subsection*{Key $L^1(\mathbb{R})$ Properties}

\noindent
\begin{tabular*}{\textwidth}{@{\extracolsep{\fill}} l l @{}}
\textbf{Property} & \textbf{Statement} \\
\hline
\textbf{Riemann--Lebesgue} & If $f \in L^1(\mathbb{R})$, then $\lim_{|\omega|\to\infty} \hat{f}(\omega) = 0$. \\
\textbf{Continuity} & If $f \in L^1(\mathbb{R})$, then $\hat{f}(\omega)$ must be a \textbf{continuous} function on $\mathbb{R}$. \\
\textbf{Inversion} & If $f, \hat{f} \in L^1(\mathbb{R})$, then $f(t) = \frac{1}{2\pi} \int_{-\infty}^{\infty} \hat{f}(\omega) e^{i\omega t} d\omega$.
\end{tabular*}
\vspace{1em}

\subsection*{Operational Rules}

\noindent
\begin{tabular*}{\textwidth}{@{\extracolsep{\fill}} l l @{}}
\textbf{Time Domain} & \textbf{Frequency Domain} \\
\hline
$\mathcal{F}\{f'(t)\}$ & $i\omega \hat{f}(\omega)$ \\
$\mathcal{F}\{t f(t)\}$ & $i \frac{d}{d\omega} \hat{f}(\omega)$ \\
$\mathcal{F}\{f(t-a)\}$ & $e^{-i\omega a} \hat{f}(\omega)$ \\
$\mathcal{F}\{f*g\}$ & $\hat{f}(\omega)\hat{g}(\omega)$ (Convolution Theorem)
\end{tabular*}

\section*{Plancherel's Theorem (Parseval's Identity for Fourier Transforms)}

Plancherel's Theorem relates the $L^2$ norm of a function in the time domain $L^2(\mathbb{R})$ to the $L^2$ norm of its Fourier Transform in the frequency domain $L^2(\mathbb{R})$. This identity is essential for evaluating integrals using Fourier transforms.

\subsection*{1. Squared $L^2$-Norm}
For a square-integrable function $f \in L^2(\mathbb{R})$, the $L^2$-norm is preserved under the Fourier Transform:
\[
  \int_{-\infty}^{\infty} |f(t)|^2 dt = \frac{1}{2\pi} \int_{-\infty}^{\infty} |\hat{f}(\xi)|^2 d\xi
\]
\emph{Note:} The factor $\frac{1}{2\pi}$ depends on the convention used for the Fourier Transform.

\subsection*{2. Polarized Plancherel Formula}
For $f, g \in L^2(\mathbb{R})$, the inner product is preserved:
\[
  \int_{-\infty}^{\infty} f(t)\overline{g(t)} dt = \frac{1}{2\pi} \int_{-\infty}^{\infty} \hat{f}(\xi)\overline{\hat{g}(\xi)} d\xi
\]
The integrals in Problem 6(b) often rely on applying these formulas, usually by identifying the integrand with either the square of a transform $|\hat{f}(\xi)|^2$ or a product of transforms $\hat{f}(\xi)\overline{\hat{g}(\xi)}$.

\end{document}